\documentclass[12pt,a4paper]{article}
\usepackage[utf8]{inputenc}
\usepackage[T1]{fontenc}
\usepackage[spanish]{babel}
\usepackage{geometry}
\usepackage{xcolor}
\usepackage{booktabs}
\usepackage{longtable}
\usepackage{array}
\usepackage{colortbl}
\usepackage{tikz}
\usetikzlibrary{shapes.geometric, arrows.meta, positioning, calc, fit, backgrounds,
                decorations.pathreplacing}
\usepackage{hyperref}
\usepackage{fancyhdr}
\usepackage{titlesec}
\usepackage{amsmath}
\usepackage{amssymb}
\usepackage{enumitem}
\usepackage{tcolorbox}
\tcbuselibrary{skins, breakable}
\usepackage{pgfplots}
\pgfplotsset{compat=1.18}

\geometry{top=2.5cm, bottom=2.5cm, left=2.8cm, right=2.8cm}
\setlength{\headheight}{15pt}

% -- Colors ------------------------------------------------------------------
\definecolor{primary}{RGB}{108,99,255}
\definecolor{secondary}{RGB}{167,139,250}
\definecolor{success}{RGB}{76,175,138}
\definecolor{errorred}{RGB}{224,92,92}
\definecolor{warnyellow}{RGB}{245,158,11}
\definecolor{blueface}{RGB}{96,165,250}
\definecolor{pinkface}{RGB}{244,114,182}

% -- Header/Footer -----------------------------------------------------------
\pagestyle{fancy}
\fancyhf{}
\fancyhead[L]{\textcolor{primary}{\textbf{CogniFace}} --- Informe del Proyecto}
\fancyhead[R]{\textcolor{gray}{\small UTN --- FECYT --- Psicologia}}
\fancyfoot[C]{\thepage}
\renewcommand{\headrulewidth}{0.4pt}

% -- Titles ------------------------------------------------------------------
\titleformat{\section}[block]{\large\bfseries\color{primary}}{\thesection}{1em}{}[\titlerule]
\titleformat{\subsection}[block]{\normalsize\bfseries\color{primary!80!black}}{\thesubsection}{1em}{}

% -- Boxes -------------------------------------------------------------------
\tcbset{
  infobox/.style={enhanced, breakable,
    colback=primary!5!white, colframe=primary!50!white,
    fonttitle=\bfseries\color{primary}, arc=3mm},
  successbox/.style={enhanced, breakable,
    colback=success!8!white, colframe=success!60!white,
    fonttitle=\bfseries\color{success!70!black}, arc=3mm},
  warnbox/.style={enhanced, breakable,
    colback=warnyellow!8!white, colframe=warnyellow!60!white,
    fonttitle=\bfseries\color{warnyellow!70!black}, arc=3mm}
}

% ============================================================================
\begin{document}

% -- Portada -----------------------------------------------------------------
\begin{titlepage}
\begin{tikzpicture}[remember picture, overlay]
  \fill[primary!10!white] (current page.north west)
    rectangle ([yshift=-5.5cm]current page.north east);
\end{tikzpicture}

\vspace*{2.5cm}
\begin{center}
  {\LARGE\bfseries\textcolor{primary}{UNIVERSIDAD TECNICA DEL NORTE}}\\[0.4cm]
  {\large Facultad de Educacion, Ciencia y Tecnologia (FECYT)}\\[0.2cm]
  {\large Carrera de Psicologia}\\[2cm]

  \begin{tcolorbox}[infobox, width=0.92\textwidth]
    \centering
    {\Huge\bfseries\textcolor{primary}{CogniFace}}\\[0.5cm]
    {\LARGE Informe del Proyecto de Investigacion}\\[0.3cm]
    {\large Hipotesis de Afinidad Facial (IAF):\\
    Ventaja de Memoria de Trabajo hacia el Sexo Opuesto}
  \end{tcolorbox}

  \vspace{1.5cm}
  \begin{tabular}{rl}
    \textbf{Area de investigacion:}  & Psicologia Cognitiva / Memoria de Trabajo \\[0.3cm]
    \textbf{Diseno experimental:}    & N-Back con estimulos faciales \\[0.3cm]
    \textbf{Hipotesis (IAF):}        &
      $\text{acc}(\text{sexo opuesto}) > \text{acc}(\text{mismo sexo})$ \\[0.3cm]
    \textbf{Plataforma:}             & \texttt{https://cogniface.web.app} \\[0.3cm]
    \textbf{Estado del proyecto:}    &
      \textcolor{success}{\textbf{Operativo --- Recoleccion activa}} \\[0.3cm]
    \textbf{Fecha:}                  & Junio 2026
  \end{tabular}
\end{center}
\end{titlepage}

\tableofcontents
\newpage

% ============================================================================
\section{Introduccion}

\subsection{Contexto Cientifico}

La \textbf{memoria de trabajo} (MT) es el sistema cognitivo responsable de mantener
y manipular informacion de manera temporal mientras se ejecutan tareas complejas
(Baddeley, 2000). La \textbf{tarea N-Back} ---propuesta originalmente por Kirchner
(1958)--- es uno de los paradigmas mas utilizados para medir la capacidad de
actualizacion de informacion en tiempo real.

En el \textbf{N-Back facial}, el participante observa una secuencia de rostros y debe
indicar cuando el rostro actual coincide con el de $N$ posiciones atras. Esta variante
es especialmente relevante en cognicion social, ya que los rostros constituyen una
clase de estimulos biologicamente privilegiada.

\subsection{Hipotesis de Afinidad Facial}

\begin{tcolorbox}[infobox, title={Hipotesis central}]
\textit{``Los participantes mostraran mayor precision en la deteccion de rostros del
sexo opuesto frente a rostros del mismo sexo en la tarea N-Back, evidenciando
una ventaja diferencial en la memoria de trabajo social para estimulos faciales
con valencia de afinidad.''}
\end{tcolorbox}

\subsection{Objetivo General}

Desarrollar y validar una plataforma web experimental para administrar la tarea
N-Back con rostros humanos y cuantificar el Indice de Afinidad Facial (IAF) en
una muestra de estudiantes universitarios de la UTN-FECYT.

% ============================================================================
\section{Metodologia}

\subsection{Diseno Experimental}

El estudio emplea un diseno \textbf{intra-sujeto} de una sola sesion. El investigador
asigna la condicion (N=1 o N=2) antes de que el participante comience.

\begin{longtable}{>{\bfseries}p{4.5cm} p{8.5cm}}
\toprule
Variable & Descripcion \\
\midrule
\endfirsthead
\toprule
Variable & Descripcion \\
\midrule
\endhead
\bottomrule
\endfoot
Variable independiente  & Sexo del rostro mostrado (masculino / femenino) \\
Variable dependiente 1  & Precision: proporcion de respuestas correctas \\
Variable dependiente 2  & Tiempo de reaccion (RT) en milisegundos \\
Variable dependiente 3  & IAF = acc(opuesto) - acc(mismo) \\
Variable de control 1   & Sexo del participante (registrado en formulario) \\
Variable de control 2   & Condicion N-Back (N=1 basico / N=2 avanzado) \\
\end{longtable}

\subsection{Participantes}

\begin{tcolorbox}[infobox, title={Criterios de inclusion}]
\begin{itemize}[leftmargin=1em]
  \item Estudiantes de la UTN-FECYT (cualquier carrera)
  \item Edad entre 18 y 35 anos
  \item Sin diagnostico reportado de deficit de atencion o trastorno de memoria
  \item Vision normal o corregida a normal
  \item Consentimiento informado verbal (protocolo de anonimato)
\end{itemize}
\end{tcolorbox}

\subsection{Estimulos}

Los estimulos consisten en \textbf{12 fotografias estandarizadas de rostros}:
6 femeninos (f01--f06) y 6 masculinos (m01--m06). Todas las imagenes presentan
expresion neutra y se precargan al inicio de la sesion para eliminar variabilidad
de latencia de red.

\begin{center}
\begin{tikzpicture}[
  face/.style={rectangle, draw=#1, fill=#1!12, rounded corners=3pt,
               minimum width=1.4cm, minimum height=1.6cm,
               align=center, font=\tiny\bfseries}
]
\foreach \i/\lbl in {0/f01, 1/f02, 2/f03, 3/f04, 4/f05, 5/f06}{
  \node[face=pinkface] at (\i*1.6, 0) {\large\textcolor{pinkface}{F}\\\lbl};
}
\foreach \i/\lbl in {0/m01, 1/m02, 2/m03, 3/m04, 4/m05, 5/m06}{
  \node[face=blueface] at (\i*1.6, -2.2) {\large\textcolor{blueface}{M}\\\lbl};
}
\node[font=\small\color{gray}, anchor=north] at (4, -3.2)
  {12 rostros unicos estandarizados (6 femeninos + 6 masculinos)};
\end{tikzpicture}
\end{center}

\subsection{Procedimiento}

\begin{center}
\begin{tikzpicture}[
  flowbox/.style={rectangle, draw=primary, fill=primary!8, rounded corners=3pt,
               minimum width=2.6cm, minimum height=0.9cm,
               font=\small\bfseries, align=center},
  arr/.style={-{Stealth[length=5pt]}, thick, color=primary!70}
]
\node[flowbox] (p0) at (0,0)   {Consentimiento};
\node[flowbox, right=1.5cm of p0] (p1) {Formulario};
\node[flowbox, right=1.5cm of p1] (p2) {Selector N};
\node[flowbox, right=1.5cm of p2] (p3) {Instrucciones};
\node[flowbox, right=1.5cm of p3] (p4) {Practica x5};
\node[flowbox, below=1cm of p4] (p5) {Experim.\\x20};
\node[flowbox, left=1.5cm of p5, fill=success!12, draw=success] (p6)
  {Resultados};
\node[flowbox, left=1.5cm of p6, fill=gray!10, draw=gray] (p7)
  {Datos\\guardados};

\draw[arr] (p0)--(p1); \draw[arr] (p1)--(p2); \draw[arr] (p2)--(p3);
\draw[arr] (p3)--(p4); \draw[arr] (p4)--(p5); \draw[arr] (p5)--(p6);
\draw[arr] (p6)--(p7);
\end{tikzpicture}
\end{center}

\subsection{Temporizado por Ensayo}

\begin{center}
\begin{tikzpicture}[scale=0.9]
\fill[gray!20, rounded corners=2pt] (0,-0.35) rectangle (10.5, 0.35);
\fill[gray!45, rounded corners=2pt] (0,-0.35) rectangle (1.5, 0.35);
\fill[primary!40, rounded corners=2pt] (1.5,-0.35) rectangle (4.3, 0.35);
\fill[secondary!40, rounded corners=2pt] (4.3,-0.35) rectangle (10.3, 0.35);

\node[font=\bfseries\footnotesize] at (0.75, 0)  {+ 500 ms};
\node[font=\bfseries\footnotesize] at (2.9,  0)  {Rostro 1000 ms};
\node[font=\bfseries\footnotesize] at (7.3,  0)  {Resp. 2000 ms};
\node[font=\small\bfseries] at (11.0, 0) {= 3500 ms};

\draw[decorate, decoration={brace, mirror, amplitude=4pt}]
  (0,-0.55) -- (1.5,-0.55) node[midway, below=4pt, font=\footnotesize]{Fijacion};
\draw[decorate, decoration={brace, mirror, amplitude=4pt}]
  (1.5,-0.55) -- (4.3,-0.55) node[midway, below=4pt, font=\footnotesize]{Estimulo};
\draw[decorate, decoration={brace, mirror, amplitude=4pt}]
  (4.3,-0.55) -- (10.3,-0.55) node[midway, below=4pt, font=\footnotesize]{Ventana resp.};
\end{tikzpicture}
\end{center}

% ============================================================================
\section{Metricas de Medicion}

\begin{longtable}{>{\bfseries}p{3.8cm} p{4.5cm} p{5cm}}
\toprule
Metrica & Formula & Interpretacion \\
\midrule
\endfirsthead
\toprule
Metrica & Formula & Interpretacion \\
\midrule
\endhead
\bottomrule
\endfoot
Aciertos (Hits) &
  Target detectado & Memoria N-Back exitosa \\[0.4em]

Omisiones (Misses) &
  Target no detectado & Fallo de actualizacion \\[0.4em]

Errores (FA) &
  Respuesta ante distractor & Inhibicion fallida \\[0.4em]

Error Emocional &
  Errores en rostros del sexo opuesto & Perturbacion por afinidad \\[0.4em]

Tiempo de Reaccion &
  $\overline{RT} = \frac{\sum RT_i}{n_{\text{hits}}}$ & Velocidad de procesamiento \\[0.4em]

Precision &
  $\text{acc} = \frac{\text{hits}+\text{CR}}{\text{total}}$ & Desempeno global \\[0.4em]

\textbf{IAF} &
  $\text{acc}_{\text{opuesto}} - \text{acc}_{\text{mismo}}$ &
  \textbf{Hipotesis de afinidad} \\
\end{longtable}

\subsection{Clasificacion de Respuestas}

\begin{center}
\renewcommand{\arraystretch}{1.5}
\begin{tabular}{|c|c|c|}
\hline
\rowcolor{primary!15}
& \textbf{RESPONDE (ESPACIO)} & \textbf{NO RESPONDE} \\
\hline
\textbf{TARGET} &
\cellcolor{success!20}\textcolor{success!70!black}{Acierto (Hit) acc=1} &
\cellcolor{errorred!20}\textcolor{errorred!70!black}{Omision (Miss) acc=0} \\
\hline
\textbf{DISTRACTOR} &
\cellcolor{warnyellow!20}\textcolor{warnyellow!70!black}{Falsa Alarma acc=0} &
\cellcolor{gray!15}{Rechazo Correcto acc=1} \\
\hline
\end{tabular}
\end{center}

% ============================================================================
\section{Calculo del IAF}

\subsection{Formula Completa}

\begin{tcolorbox}[infobox, title={Indice de Afinidad Facial por sexo del participante}]
\textbf{Participante masculino:}
\begin{equation}
  \text{IAF} = \frac{\#\text{correctos en F}}{\#\text{ensayos F}} -
               \frac{\#\text{correctos en M}}{\#\text{ensayos M}}
\end{equation}

\textbf{Participante femenino:}
\begin{equation}
  \text{IAF} = \frac{\#\text{correctos en M}}{\#\text{ensayos M}} -
               \frac{\#\text{correctos en F}}{\#\text{ensayos F}}
\end{equation}

\medskip
Donde $F$ = ensayos con rostros femeninos y $M$ = ensayos con rostros masculinos
(solo ensayos reales, sin practica).

\bigskip
\centering
\begin{tabular}{cll}
\toprule
\textbf{Valor IAF} & \textbf{Interpretacion} & \textbf{Hipotesis} \\
\midrule
$> 0$           & Ventaja hacia sexo opuesto & \textcolor{success}{\bfseries Apoyada} \\
$= 0$           & Sin diferencia             & Nula \\
$< 0$           & Ventaja hacia mismo sexo   & \textcolor{errorred}{\bfseries No apoyada} \\
\bottomrule
\end{tabular}
\end{tcolorbox}

\subsection{Ejemplo Ilustrativo}

\begin{tcolorbox}[successbox, title={Ejemplo: Participante Masculino}]
\begin{tabular}{lll}
Ensayos con rostros femeninos: & 8 de 10 correctos & $\text{acc}(F) = 0.80$ \\
Ensayos con rostros masculinos: & 6 de 10 correctos & $\text{acc}(M) = 0.60$ \\
\end{tabular}\\[0.5em]
$\text{IAF} = 0.80 - 0.60 = \mathbf{+0.200}$\\[0.3em]
IAF positivo $\Rightarrow$ hipotesis de afinidad \textbf{apoyada} para este participante.
\end{tcolorbox}

\subsection{IAF esperado por grupo (simulacion hipotetica)}

\begin{center}
\begin{tikzpicture}
\begin{axis}[
  width=11cm, height=6cm,
  ylabel={IAF promedio},
  xlabel={Grupo y condicion},
  ymin=-0.1, ymax=0.45,
  xtick={1,2,3,4},
  xticklabels={Hombres N1, Mujeres N1, Hombres N2, Mujeres N2},
  ymajorgrids=true,
  grid style={dashed, gray!40},
  ybar, bar width=1cm,
  legend pos=north east,
  title={Distribucion IAF esperada (hipotesis)}
]
\addplot[fill=blueface!60, draw=blueface!80] coordinates {
  (1,0.18) (2,0.00) (3,0.12) (4,0.00)
};
\addplot[fill=pinkface!60, draw=pinkface!80] coordinates {
  (1,0.00) (2,0.22) (3,0.00) (4,0.15)
};
\draw[dashed, thick, errorred] (axis cs:0.3,0) -- (axis cs:4.7,0);
\addlegendentry{Hombres}
\addlegendentry{Mujeres}
\end{axis}
\end{tikzpicture}
\end{center}

% ============================================================================
\section{Estructura de la Sesion Experimental}

\subsection{Fase de Practica (5 ensayos)}

Durante los 5 ensayos de practica el participante recibe retroalimentacion inmediata:

\begin{center}
\begin{tabular}{cll}
\toprule
Respuesta & Feedback visual & Color \\
\midrule
Correcto (Hit o CR) & ``\textbf{Correcto!}'' & \textcolor{success}{Verde} \\
Incorrecto (Miss o FA) & ``\textbf{Incorrecto}'' & \textcolor{errorred}{Rojo} \\
\bottomrule
\end{tabular}
\end{center}

Los ensayos de practica \textbf{no} se incluyen en el analisis de datos.

\subsection{Fase Experimental (20 ensayos)}

Los 20 ensayos reales se presentan sin retroalimentacion. La distribucion
garantiza exactamente 6 targets por bloque, 3 de cada sexo:

\begin{center}
\begin{tikzpicture}
\begin{axis}[
  width=9cm, height=5.5cm,
  ybar stacked,
  bar width=1.8cm,
  xtick={1,2},
  xticklabels={N-Back 1, N-Back 2},
  ylabel={Numero de ensayos},
  ymin=0, ymax=22,
  ymajorgrids=true,
  grid style={dashed, gray!30},
  legend style={at={(0.98,0.98)}, anchor=north east, font=\small},
  title={Composicion de ensayos por bloque}
]
\addplot[fill=pinkface!60, draw=pinkface] coordinates {(1,3) (2,3)};
\addplot[fill=blueface!60, draw=blueface] coordinates {(1,3) (2,3)};
\addplot[fill=gray!30, draw=gray!60] coordinates {(1,14) (2,14)};
\addlegendentry{Target femenino (3)}
\addlegendentry{Target masculino (3)}
\addlegendentry{Distractor (14)}
\end{axis}
\end{tikzpicture}
\end{center}

Ratio de targets: $6/20 = 30\%$ por bloque, equilibrado por sexo de rostro.

% ============================================================================
\section{Balance del Proyecto}

\subsection{Estado de Implementacion}

\begin{longtable}{>{\bfseries\small}p{0.5cm} p{6.2cm} c p{4.5cm}}
\toprule
\# & Requisito & Estado & Notas \\
\midrule
\endfirsthead
\toprule
\# & Requisito & Estado & Notas \\
\midrule
\endhead
\bottomrule
\endfoot
1 & Dark/Light mode en todas las interfaces
  & \textcolor{success}{\bfseries Completo}
  & ThemeToggle en 7 pantallas \\

2 & Pantalla bienvenida UTN/FECYT/Psicologia
  & \textcolor{success}{\bfseries Completo}
  & WelcomeScreen con branding institucional \\

3 & Descripcion de la tarea cognitiva
  & \textcolor{success}{\bfseries Completo}
  & Texto cientifico en WelcomeScreen \\

4 & Formulario sexo, ID participante, edad
  & \textcolor{success}{\bfseries Completo}
  & Validacion, botones masculino/femenino \\

5 & Selector N-Back 1 o N-Back 2
  & \textcolor{success}{\bfseries Completo}
  & BlockSelector con tarjetas descriptivas \\

6 & Instrucciones por bloque elegido
  & \textcolor{success}{\bfseries Completo}
  & Timing, ejemplos visuales, reglas \\

7 & 5 ensayos de practica con feedback
  & \textcolor{success}{\bfseries Completo}
  & Retroalimentacion inmediata correcto/incorrecto \\

8 & 20 ensayos experimentales sin feedback
  & \textcolor{success}{\bfseries Completo}
  & ExperimentBlock sin overlay \\

9 & Resultados: Aciertos, Omisiones, TR, Error, Error Emocional
  & \textcolor{success}{\bfseries Completo}
  & ResultsScreen con IAF verdict \\

10 & Exportacion Excel desde panel admin
  & \textcolor{success}{\bfseries Completo}
  & SheetJS, hojas Participantes y Ensayos \\

11 & Deploy Firebase Hosting con CI/CD
  & \textcolor{success}{\bfseries Completo}
  & GitHub Actions, cogniface.web.app \\

12 & README orientado al proyecto
  & \textcolor{success}{\bfseries Completo}
  & Branding UTN/FECYT, flujo, metricas \\
\end{longtable}

\begin{tcolorbox}[successbox, title={Resultado del Balance}]
\textbf{12 de 12 requisitos implementados y verificados.}\\[0.3em]
Build de produccion: \texttt{0 errores --- 667 modulos --- 1.08 s}\\
URL activa: \texttt{https://cogniface.web.app}
\end{tcolorbox}

% ============================================================================
\section{Protocolo para el Investigador}

\subsection{Antes de la Sesion}
\begin{enumerate}[leftmargin=1.5em]
  \item Abrir \texttt{https://cogniface.web.app} en el equipo de laboratorio
  \item Asignar condicion (N=1 o N=2) al participante segun protocolo
  \item Entregar ID de participante (ej. \texttt{EST\_2026\_042})
  \item Verificar que el participante comprende la tarea
\end{enumerate}

\subsection{Despues de la Sesion}
\begin{enumerate}[leftmargin=1.5em]
  \item Acceder a \texttt{https://cogniface.web.app/admin} con credenciales
  \item Verificar que el registro aparece como ``Completado'' en la tabla
  \item Revisar el IAF calculado automaticamente
  \item Exportar Excel actualizado con todos los participantes
\end{enumerate}

\begin{tcolorbox}[warnbox, title={Seguridad del acceso}]
Las credenciales del panel admin son exclusivas del investigador.
No compartir con participantes.
\end{tcolorbox}

% ============================================================================
\section{Interpretacion de Resultados}

\begin{longtable}{>{\centering}p{2.5cm} p{5cm} p{5cm}}
\toprule
\textbf{Rango IAF} & \textbf{Interpretacion} & \textbf{Recomendacion} \\
\midrule
$> +0.10$            & Fuerte afinidad hacia sexo opuesto & Hipotesis claramente apoyada \\
$[+0.01, +0.10]$     & Afinidad leve o moderada           & Hipotesis apoyada debilmente \\
$[-0.01, +0.01]$     & Sin diferencia significativa        & Resultado nulo \\
$< -0.01$            & Ventaja hacia mismo sexo            & Hipotesis no apoyada \\
\bottomrule
\end{longtable}

\subsection{Error Emocional}

El Error Emocional cuantifica los errores en rostros del sexo opuesto:

\begin{equation}
  \text{Error Emocional} =
  \left|\left\{\,i : \text{face\_gender}_i = \text{opuesto}
  \;\land\; (\text{miss}_i \lor \text{FA}_i)\right\}\right|
\end{equation}

Un valor alto puede indicar interferencia cognitiva por la valencia de afinidad.

% ============================================================================
\section{Consideraciones Eticas}

\begin{tcolorbox}[infobox, title={Protocolo de privacidad}]
\begin{itemize}[leftmargin=1em]
  \item Participantes registrados con codigo anonimo (sin nombre ni correo)
  \item Informacion: sexo, edad, codigo asignado
  \item Datos en Firestore con acceso restringido al investigador autenticado
  \item No se comparten datos con terceros
  \item Participantes pueden abandonar en cualquier momento
  \item Consentimiento verbal e implicito al iniciar la tarea
\end{itemize}
\end{tcolorbox}

% ============================================================================
\section{Limitaciones}

\begin{enumerate}[leftmargin=1.5em]
  \item \textbf{Muestra de conveniencia:} Estudiantes UTN-FECYT, limita generalizacion.
  \item \textbf{Condicion unica:} Un solo bloque por participante (N=1 o N=2), sin
        contrabalanceo intra-sujeto.
  \item \textbf{Identidad de rostros:} Pueden ser conocidos por algunos participantes.
  \item \textbf{Sin medidas complementarias:} Sin screening de MT ni ansiedad social.
  \item \textbf{Analisis inferencial externo:} El IAF descriptivo calculado requiere
        pruebas de significancia en SPSS, R o Python para publicacion.
\end{enumerate}

% ============================================================================
\section*{Conclusiones}

CogniFace provee una plataforma completa, reproducible y de costo cero para
administrar el experimento N-Back facial en contexto universitario. La implementacion
cubre el ciclo completo: reclutamiento, administracion estandarizada de la tarea,
calculo automatico del IAF y exportacion de datos para analisis inferencial.

\begin{tcolorbox}[successbox, title={Proximos pasos para el equipo investigador}]
\begin{enumerate}[leftmargin=1.5em]
  \item Definir tamano muestral objetivo (recomendado: $n \geq 30$ por sexo
        $\times$ condicion)
  \item Establecer protocolo de aplicacion estandarizado en laboratorio
  \item Iniciar recoleccion de datos y monitorear el panel admin
  \item Exportar Excel al completar la muestra y analizar en software estadistico
  \item Redactar resultados y discusion para informe final de investigacion
\end{enumerate}
\end{tcolorbox}

\vspace{2em}
\noindent\rule{\textwidth}{0.4pt}\\[0.3em]
{\small CogniFace v1.0 --- Universidad Tecnica del Norte,
FECYT, Carrera de Psicologia --- Junio 2026}

\end{document}
